\def\level{BTS}  % Classe à droite
\def\course{CIEL 2}
\def\chaptername{Lois continues} % Titre au centre
\def\docpart{Loi uniforme~--~Exercices} % Partie du cours
\def\docname{C2251}        % Identifiant à gauche

\pagestyle{cours_FP}
\makeheaderBTS{} % Affiche le bandeau

\begin{multicols}{2}

	\begin{exercice_cours}[\quad]{}
Soit $X$ une variable aléatoire suivant une loi uniforme sur $[-3~;~4]$.

Tracer la densité de probabilité de $X$ et le domaine dont l'aire correspond à \[P(-1\leq X\leq 2)\]
	\end{exercice_cours}
	\begin{exercice_cours}[\quad]{}

Soit $X$ une variable aléatoire suivant une loi uniforme sur $[1~;~7]$.

Déterminer les probabilités suivantes :

\begin{itemize}
	\item $P(4\leq X\leq 5)$
	\item $P(X=2)$
	\item $P(X> 5)$
\end{itemize}

		\end{exercice_cours}
\begin{exercice_cours}[\quad]{}

Soit $X$ une variable aléatoire suivant une loi uniforme sur $[0~;~10]$.

Déterminer les probabilités suivantes :

\begin{itemize}
	\item $P(X<6)$
	\item $P(3\leq X\leq 7)$
	\item $P(X\geq 7)$
\end{itemize}

		\end{exercice_cours}

		\begin{exercice_cours}[\quad]{}

		$X$ est une variable aléatoire qui suit la loi uniforme sur $\ff{-2}{2}$.

		\begin{itemize}
			\item Calculer $P(X<1)$ et $P(X\geq 0,5)$
			\item Calculer $P_{X>0}(X<1)$
			\item Donner l'espérance de $X$
		\end{itemize}
	\end{exercice_cours}

	
\begin{exercice_cours}[\quad]{}


Une entreprise vient d'envoyer un mail à un client à 13h. Vu le contenu du message, il est certain que le client va répondre par mail et que ce mail peut arriver n'importe quand entre 14h et 17h.

On admet que la durée d'attente du mail renvoyé par le client est modélisée par la variable aléatoire $T$ qui suit la loi uniforme sur l'intervalle $[14~;~17]$.

\begin{enumerate}
	\item Déterminer la fonction densité de probabilité de T. 
	\item Quelle est la probabilité que le mail arrive à 15h ?
	\item Quelle est la probabilité que le mail arrive entre 14h et 15h ?
	\item Quelle est la probabilité que le mail arrive entre 16h et 17h ?
	\item Quelle est la probabilité que le mail arrive après 16h30 ?
\end{enumerate} 

		\end{exercice_cours}
\begin{exercice_cours}[\quad]{}

Suite à un problème de réseau, un client contacte le service après-vente de son opérateur.

Un conseiller l'informe qu'un technicien le rappellera pour une intervention à distance entre 14h et 15h.

Sachant que ce technicien appelle de manière aléatoire sur le créneau donné, on souhaite calculer la probabilité que le client patiente entre $15$ et $40$ minutes.


On désigne par $T$ la variable aléatoire continue qui donne le temps d’attente en minutes.

\begin{enumerate}
	\item Donner la fonction densité de probabilité associée à $X$.
	\item Répondre au problème posé.
	\item Quelle est la probabilité que le client patiente au moins $10$ minutes ?
	\item Calculer la probabilité que le client patiente entre $15$ et $40$ minutes sachant qu'il a déjà attendu $10$ minutes.
	\item Quel est le temps d'attente moyen lorsque ce genre d'interventions sont proposées un grand nombre de fois ?
\end{enumerate}

	\end{exercice_cours}

	

	\vfill
\end{multicols}

